\section{Zigzags and Tightening}

\subsection{Zigzags in a bisupport category \catcw}

For any $n \ge 1$, a zigzag of length $n$ in bisupport category \catcw  
is a sequence of morphisms  $F=\zigzagTuple{f}{n}$ of \catcw such that for some 
objects $x_0,..x_n$ and $y_1,...y_n$ of \catcw we have
$$
\arraycolsep=7pt
\zigzag{x}{y}{f}{n}
$$
in \catcw. 

Within a zigzag of length $n$ we find a zig for each \foreachi[n-1], like so
$$
\begin{array} {c c c}
&\Ob{y}{i}&\\[0.8cm]
\Ob{x}{i}&&\Obs{x}{i}
\end{array}
\begin{arrows}
\ncarr{xi}{yi}   \alabel{f_i}[0.45][0.4]
\ncarr {xis}{yi} \blabel{f'_i}[0.45][0.4]
\end{arrows}
$$

We will say that this i'th zig is \textit{tight} iff
$\rng{f_i}=\rng{f'_i}$.

We also find a  zag for each $i$, $2 \leq i \leq n$, like so
%\arreset}
$$
\setlength{\arrnodesepA}{3pt}
\begin{array} {c c c}
\Obp{y}{i}&&\Ob{y}{i} \\[0.8cm]
&\Ob{x}{i}&
\end{array}
\begin{arrows}
\ncarr{xi}{yip}  \alabel{f'_i}[0.45][0.4]
\ncarr {xi}{yi}  \blabel{f_i}[0.45][0.4]
\end{arrows}
$$

and we will say that the i'th zag is tight iff
$\rst{f'_i}=\rst{f_i}$.

A zigzag we will say is a \textit{tight zigzag} iff its every zig and its every zag are tight. 

If $F=\zigzagTuple{f}{n}$ and 
$G=\zigzagTuple{g}{m}$ are zigzags such that the codomain of $f_n$
is the domain of $g_1$ so that the morphism $f_n \circ g_1$ is defined in \catcw then the zigzags $F$ and $G$ can be composed to give a zigzag 
$F::G=\tuple{f_1,f'_1,...f'_{n-1},f_n \circ g_1,g'_1, g_2,...g_m}$ 
of length $n+m-1$ as here:
$$
\arraycolsep=7pt
\left( \zigzag{w}{x}{f}{n}\right) \cnc \left(\zigzag{y}{z}{g}{m}\right)
$$

$$
=_{def} \  \zigzagfgcomposite
$$

We will say that $F \cnc G$ is the concatenation of $F$ with $G$.
Concatenation is associative because composition in the underlying category
is and there is a category of zigzags but the category of zigzags is of no particular interest to us, rather we are interested in establishing a category of tight zigzags and this requires some work because generally the concatenation of two tight zigzags is not tight.  

Support and cosupport can be defined for zigzags from the support and cosupport in the underlying bisupport category. 
If $F$ is the zigzag  
given by tuple $\zigzagTuple{f}{n}$ then define $\rst{F}$ to be 
the singleton zigzag $\tuple{\rst{f_1}}$ and define
$\rng{F}$ to be the singleton tuple $\tuple{\rng{f_n}}$.
With support and cosupport so defined it is easy to see  that there is  a bisupport category of zigzags.  

To save some ink, if $f: a \morph b$ 
in the underlying bisupport category \catcw 
and if $G:b \morph c$ is a zigzag
over \catcw then we will write $f \cnc G$ to mean the concatenation
of the singleton zigzag $\tuple{f}$ with the zigzag $G$, i.e. we will write
$f \cnc G$ in place of $\tuple{f} \cnc G$. Likewise when $F$ is a zigzag and 
$g$ is a morphism in the underlying category then we will write
$F \cnc g$ in place of $F \cnc \tuple{g}$.

\subsection{Comparison of Zigzags}

Say that zigzags are comparable if they are of the same length
and have identical object sequences.
If $F=\zigzagTuple{f}{n}$ and 
$G=\zigzagTuple{g}{n}$
 are comparable zigzags then define $F \leq G$
by $F \leq G$ iff for each $i$, $0 \leq i \leq n$
$f_i \leq g_i$ and for each $i$, $0 \leq i \leq n-1$ $f'_i \leq g'_i$.

\subsection{Tightening}


\NewDocumentCommand{\tightlead}{o}
{%
  f_1,f'_1,\ldots,f_{n-2},f'_{n-2},
  f_{n-1}\circ
  \rng{
    \IfValueTF{#1}
      {\rst{f_n\circ #1}}
      {\rst{f_n}}
    \circ f'_{n-1}
  }%
}

\NewDocumentCommand{\leftToRightTightening}{o}
{%
\tuple{t_1,t'_1,\ldots,t'_{n-2},t_{n-1},\;
    \IfValueTF{#1}
      {\rst{f_n\circ #1}}
      {\rst{f_n}}
    \circ f'_{n-1} \circ \rng{t_{n-1}},\;
    \rst{f'_{n-1} \circ \rng{t_{n-1}}}
    \circ
    \IfValueTF{#1}
      {f_n\circ #1}
      {f_n}
    }
}

If $F$ is a zigzag in a poset-enriched bisupport category \catcw, 
a \textit{tightening} of $F$ is a tight zigzag
$T \leq F$ such that for every tight zigzag $Z$ with $Z \leq F$,
we have $Z \leq T$.

Because $\leq$ is a partial order, any such tightening of $F$ is unique.

We will show by induction on length that every zigzag $F$ has a tightening and we will then write $T(F)$ for the tightening of $F$.
If the zigzag is of length $1$ then it has no zigs and no zags so is already tight. The inductive step is given by the next lemma: 

{  %lemma leftToRightTighteningSubLemma

\begin{lemma}
\llabel{leftToRightTighteningSubLemma}
Suppose  $F=\zigzagTuple{f}{n}$ is a zigzag in range category \catcw.
If  $\zigzagLeadingTuple{t}{n}$ is a tightening of the zigzag
\begin{equation}
\label{inductionPrecursor}
\tuple{\tightlead}
\end{equation}
then  zigzag
\begin{equation}
\label{requiredTightening}\leftToRightTightening
\end{equation}
is a tightening of $F$. 
\end{lemma}
\begin{proof}
We have to show first that (\ref{requiredTightening}) is tight. 

For this we need the following identities:
\begin{align}
\label{tightnessIdentityOne}
\rng{t_{n-1}} &= \rng{\rst{f_n} \circ f'_{n-1} \circ \rng{t_{n-1}}}  \\[0.1cm]
\label{tightnessIdentityTwo}
\rst{\rst{f_n} \circ f'_{n-1} \circ \rng{t_{n-1}}}&=\rst{\rst{f'_{n-1} \circ \rng{t_{n-1}}} \circ f_n}
\end{align}

To show that (\ref{tightnessIdentityOne}) holds first note that 
because from the definition of $\zigzagLeadingTuple{t}{n}$ as a tightening
it follows in particular that
\[
t_{n-1} \leq f_{n-1} \circ \rng{\rst{f_n} \circ f'_{n-1}}
\]
and that therefore that
\[
\rng{t_{n-1}} \leq \rng{f_{n-1} \circ \rng{\rst{f_n} \circ f'_{n-1}}}
\]
which rearranges using C.3 to
\[
\rng{t_{n-1}} \leq \rng{f_{n-1}} \circ \rng{\rst{f_n} \circ f'_{n-1}}
\]
from which we have by lemma \lref{idempotentCompositionAndOrdering}(ii)
that 
\[
\rng{t_{n-1}} \leq \rng{\rst{f_n} \circ f'_{n-1}}.
\]
From this inequality it follows that
\begin{align*}
\rng{t_{n-1}} 
           &= \rng{t_{n-1}} \circ \rng{\rst{f_n} \circ f'_{n-1}}
           && \mbox{by lemma \lref{idempotentCompositionAndOrdering}(viii),}\\
           &= \rng{\rst{f_n} \circ f'_{n-1}} \circ \rng{t_{n-1}}
           && \mbox{by C.2,}\\
           &= \rng{\rst{f_n} \circ f'_{n-1} \circ \rng{t_{n-1}}}
           &&\mbox{by C.3, as required.}
\end{align*}

Identity (\ref{tightnessIdentityTwo}) holds \highlight{CHECK} because both sides  simplify, using lemma \lref{bisupportLemma}(i), and axiom S.2, 
to $\rst{f'_{n-1} \circ \rng{t_{n-1}}} \circ \rst{f_n}$.

We can now show that (\ref{requiredTightening})
is the tightening of $F$ by showing that 
it dominates any tight zigzag $Z$ such that
$Z \leq F$ 
Assume such a $Z=\zigzagTuple{z}{n}$ and assume that

\[
\zigzagTuple{z}{n} \leq \zigzagTuple{f}{n}
\]
we need to show that 
\begin{equation}
\label{targetInequality}
Z \leq \leftToRightTightening.
\end{equation} 

First we note that  
\[
 z_{n-1} \leq f_{n-1} \circ 
       \rng{
                \rst{f_n} \circ f'_{n-1}
                     }
\]
because 
\begin{align*}
z_{n-1} &= z_{n-1} \circ \rng{\rst{z_n} \circ z'_{n-1}} && \mbox{trivially, since $Z$ is tight}\\
&\leq  f_{n-1} \circ 
       \rng{
                \rst{f_n} \circ f'_{n-1}
                     } &&\mbox{\parbox{3cm}{since $Z \leq F$.}}
\end{align*}

Thus since also $Z \leq F$ then we have that
\[
\zigzagLeadingTuple{z}{n} \leq \tuple{\tightlead}
\] 
and from this and the definition of $\zigzagLeadingTuple{t}{n}$
as a tightening of the rhs of this inequality we establish that
\begin{equation}
\label{inequalityOne}
\zigzagLeadingTuple{z}{n} \leq \zigzagLeadingTuple{t}{n}
\end{equation}
 
To establish the target inequality (\ref{targetInequality})
we need 
in addition to show that
 \begin{equation}
 \label{inequalityTwo}
 z'_{n-1} \leq \rst{f_n} \circ f'_{n-1} \circ \rng{t_{n-1}}
 \end{equation}
 and
  \begin{equation}
 \label{inequalityThree}
 z_n \leq \rst{f'_{n-1} \circ \rng{t_{n-1}}} \circ f_n
 \end{equation}

The first of these holds because
\begin{align*}
z'_{n-1} &= \rst{z'_{n-1}} \circ z'_{n-1}  \circ \rng{z'_{n-1}} 
                && \mbox{trivially, by axioms S.1 and C.1,}\\
         &= \rst{z_n} \circ z'_{n-1}  \circ \rng{z_{n-1}} 
         &&\mbox{because $Z$ is tight,}\\
         &\leq \rst{f_n} \circ f'_{n-1}  \circ \rng{z_{n-1}}
         &&\mbox{because $Z \leq F$ and support and cosupport are monotone,}\\
         &\leq \rst{f_n} \circ f'_{n-1}  \circ \rng{t_{n-1}}
         &&\mbox{because $z_{n-1} \leq t_{n-1}$ is established
                               above in \workt{(\ref{inequalityOne})}.}
\end{align*}

The second  because
\begin{align*}
 z_n &= \rst{z_n} \circ z_n
         &&\mbox{by axiom S.1,}\\        
     &=  \rst{z'_{n-1}} \circ z_n 
         &&\mbox{since $Z$ is tight,}\\
     &\leq  \rst{\rst{f_n} \circ\ f'_{n-1} \circ \rng{t_{n-1}}} \circ f_n
         &&\mbox{by (\ref{inequalityTwo}) and since $z_n \leq f_n$,} \\
         &= \rst{f'_{n-1} \circ \rng{t_{n-1}}} \circ f_n
         && \mbox{by axioms S.3, S.2 and S.1.}
\end{align*}
\end{proof}
} %end lemma leftToRightTighteningSubLemma
We have established that every zigzag $F$ has a unique tightening. 
We shall write $T(F)$ for the tightening of $F$. 


Lemma \lref{leftToRightTighteningSubLemma} describes how a zigzag can be tighten from left to right. Duality gives us the following right to left tightening lemma:

\NewDocumentCommand{\tightTail}{o}
{%
  \IfValueTF{#1}
    {\rst{f'_1 \circ \rng{#1\circ f_1}}}
    {\rst{f'_1 \circ \rng{f_1}}}
  \circ f_2,\;
  f'_2,\ldots,f_{n-2},f'_{n-1},f_n
}
\NewDocumentCommand{\rightToLeftTightening}{o}
{%
\tuple{
  \IfValueTF{#1}
    {#1\circ f_1 \circ \rng{\rst{t_2}\circ f'_1}}
    {f_1\circ \rng{\rst{t_2}\circ f'_1}},
  \;
  \rst{t_2}\circ f'_1\circ
  \IfValueTF{#1}
    {\rng{#1\circ f_1}}
    {\rng{f_1}},
  \;
  t_2,\ldots,t_{n-1},t'_{n-1},t_n
  }
}
\begin{lemma}
\llabel{rightToLeftTighteningSubLemma}
Suppose  $F=\zigzagTuple{f}{n}$ is a zigzag in range category \catcw.
If  $\zigzagTrailingTuple{t}{n}$ is a tightening of the zigzag
\begin{equation}
\label{inductionRLprecursor}
\tuple{\tightTail}
\end{equation}
then  zigzag
\begin{equation}
\label{requiredRLTightening}\rightToLeftTightening
\end{equation}
is a tightening of $F$. 
\end{lemma}

\begin{lemma}
\llabel{endpointLemma}
If $F=\zigzagTuple{f}{n}$ and $T(F)=\zigzagTuple{t}{n}$ then
\begin{enumerate}[(i)]
\item $t_1 = f_1 \circ \rng{t_1} = f_1 \circ \rng{t'_1}$,
\item $t_n = \rst{t_{n-1}} \circ f_n = \rst{t'_{n-1}} \circ f_n$.
\end{enumerate}
\end{lemma}
\begin{proof}
~ % go to new line
\begin{enumerate}[(i)]
\item
By lemma \lref{rightToLeftTighteningSubLemma}
$t_1=f_1 \circ \rng{x}$
for a certain morphism $x$.
From this starting point the required 
$t_1 = f_1 \circ \rng{t_1} = f_1 \circ \rng{t'_1}$
is shown as follows
\begin{align*}
t_1 &= f_1 \circ \rng{x} 
        &&\mbox{the starting point,} \\
    &= f_1 \circ \rng{f_1} \circ \rng{x}
        && \mbox{by axiom C.1,}  \\
    &= f_1 \circ \rng{f_1 \circ \rng{x}}
        && \mbox{by  axiom C.3,} \\
    &= f_1 \circ \rng{t_1}
        && \mbox{using $t_1=f_1 \circ \rng{x}$ again,} \\
    &= f_1 \circ \rng{t'_1}
        && \mbox{since $\zigzagTuple{t}{n}$ is tight.}
\end{align*}
\item
This follows by duality. The given zigzags $F$ and $T(F)$ give rise to zigzags
in the zigzag category over $\catc^{op}$ and the definition of tightening is such that one zigzag is the tightening over another in zigzags over the opposite category iff their duals are in the same relationship. The dual to (i) is (ii). Because (i) holds in zigzags over the opposite category then (ii) holds in zigzags over \catc.  
\end{enumerate}
\end{proof}
\begin{lemma}
\llabel{tighteningCommutesWithConcatenatingAnIdempotentSublemma}
For any zigzag $F$ and for any comparable tight zigzag $H$
\begin{enumerate}[(i)]
\item if $\tuple{i} \cnc F$ is defined for some morphism $i$ then
if  $H \leq \tuple{\rng{i}} \cnc F$
then $H \leq  \tuple{\rng{i}} \cnc \tighten{F}$.
\item if $F\cnc\tuple{j}$ is defined for some morphism $j$ then
if  $H \leq F \cnc \tuple{\rst{j}}$
then $H \leq  \tighten{F} \cnc \tuple{\rst{j}}$.
\end{enumerate}
\end{lemma}
\begin{proof}
~
Suppose $\tighten{F}$ is the zigzag $\zigzagTuple{t}{n}$
and suppose $H$ to be the zigzag $\zigzagTuple{h}{n}$.
\begin{enumerate}[(i)]
\item % (i)
Because $\tuple{\rng{i}} \cnc F \leq F$ (lemma \workt{from poset-enriched section})then we have that
$H \leq F$. Therefore since $H$ is tight and from the definition
of tightness we have that
$H \leq \tighten{F}$.

Because of this and since $\tighten{F}$ differs from
$\tuple{\rng{i}} \cnc \tighten{F}$ only in the first element
it remains only to show that
$h_1 \leq \rng{i} \circ t_1$ which we can do as follows:
\begin{align*}
h_1 &=  h_1 \circ \rng{h_1} 
             && \mbox{by axiom C.1}\\
    &\leq h_1\circ \rng{t'_1} 
             && \mbox{since $h_1 = h'_1 \leq t'_1$ and composition is monotone} \\
    &\leq \rng{i} \circ f_1 \circ \rng{t'_1} 
             && \mbox{since $h_1 \leq \rng{i} \circ f_1$} \\
    &= \rng{i \circ t_1}
            && \mbox{by \lref{endpointLemma} (i), as required.} 
\end{align*}


\item % (ii) 
Follows by duality since it is the dual of (i).
\end{enumerate}
\end{proof}
\begin{lemma}
\llabel{tighteningCommutesWithConcatenatingAnIdempotent}
For any zigzag $F$
\begin{enumerate}[(i)]
\item
if $\tuple{i} \cnc F$ is defined for some morphism $i$ then
$\tighten{\rng{i} \cnc F}
         = \tighten{\rng{i} \cnc \tighten{F}}$,
\item 
if $F \cnc \tuple{j}$ is defined for some morphism $j$ then
$\tighten{F \cnc \rst{j}}
         = \tighten{\tighten{F} \cnc \rst{j}}$.
\end{enumerate}
\end{lemma}
\begin{proof}
~
\begin{enumerate}[(i)]
\item %(i)
To show that the tightening of $\rng{i} \cnc F$
is $\tighten{\rng{i}\cnc \tighten{F}}$ it suffices to show that if
$H$ is a tight zigzag and $H \leq \rng{i} \cnc F$ then
$H \leq \tighten{\rng{i}\cnc \tighten{F}}$. So assume such an $H$.
Then by \lref{tighteningCommutesWithConcatenatingAnIdempotentSublemma} (i), $H \leq \rng{i} \cnc \tighten{F}$.
But then, because $H$ is tight,
$H \leq \tighten{\rng{i}\cnc \tighten{F}}$.
\item %(ii)
Follows by duality since it is the dual of (i).
\end{enumerate}
\end{proof}

\begin{lemma}
\llabel{alpha}
if $F=\zigzagTuple{f}{n}$ is a tight zigzag then
\begin{enumerate}[(i)]
\item if $\tuple{i} \cnc F$ is defined, 
for some morphism $i$, 
then the first element of $\tighten{\tuple{i} \cnc F}$ 
is $i \circ f_1$, 
\item if $F \cnc \tuple{j}$ is defined, 
for some morphism $j$, 
then the final element of $\tighten{F \cnc \tuple{j}}$ 
is $f_n \circ j$. 
\end{enumerate}
\end{lemma}
\begin{proof}
~
\begin{enumerate}[(i)]
\item %(i)
We use the rightToLeftTighteningSublemma to tighten 
$\tuple{i} \cnc F$ and accordingly 
we let $\zigzagTrailingTuple{t}{n}$
be the tightening of
\[
\tuple{\tightTail[i]}
\]
so that 
\begin{equation}
\label{alphaInitialSchematic}
\rightToLeftTightening[i]
\end{equation}
is the tightening of $\tuple{i}\cnc F$.

\newcommand{\inductivets}
{\rst{ f'_1 \circ \rng{i \circ f_1}}\circ f_2}

We assume the inductive hypothesis that
\[
t_2 =  \inductivets
\]
Substituting into the initial element of
(\ref{alphaInitialSchematic})
the initial element of $\tighten{i\cnc F}$
is
\[
i\circ f_1 \circ \rng{\rst{\inductivets}\circ f'_1}
\]
This can be simplified as follows:

\begin{align*}
i\circ f_1\circ
\rng{\rst{\inductivets}\circ f'_1}
&=
i\circ f_1\circ
\rng{
 \rst{f'_1\circ\rng{i\circ f_1}}
 \circ \rst{f_2}
 \circ f'_1
}
&&\mbox{by R.3,} \\
&=
i\circ f_1\circ
\rng{
 \rst{f'_1\circ\rng{i\circ f_1}}
 \circ f'_1
}
&&\mbox{since $F$ is tight and by axiom S.1,} \\
&=
i\circ f_1\circ
\rng{f'_1}
&&\mbox{\parbox[t]{5cm}{by lemma \lref{bisupportLemma} (iii), 
with $i \circ f_1$ as $f$ and $f'_1$ as $g$} }\\
&=
i\circ f_1
&&\mbox{by axiom C.1, again and as required.}
\end{align*}
\item %(ii)
Follows by duality since it is the dual of (i).
\end{enumerate}
\end{proof}

\begin{lemma}
\llabel{beta}
If $F=\zigzagTuple{f}{n}$ is a any zigzag then
\begin{enumerate}[(i)]
\item if $\tuple{i} \cnc F$ is defined, 
for some morphism $i$,  
then 
$$\rst{\tighten{\bigl\langle \rng{i} \bigr\rangle \cnc F}}
                = \tuple{\rng{i}} \circ \rst{\tighten{F}}$$, 
\item if $F \cnc \tuple{j}$ is defined, 
for some morphism $j$, then 
$$\rng{\tighten{F \cnc \tuple{\rst{j}}}} 
                = \rng{\tighten{F}} \circ \tuple{\rst{j}}$$. 
\end{enumerate}
\end{lemma}
\begin{proof}
Suppose $\zigzagTuple{t}{n}$ is the tightening of $F$.
\begin{enumerate}[(i)]
\item %(i)
$\tighten{\tuple{\rng{i}} \cnc F}$ is
$\tighten{\tuple{\rng{i}} \cnc \tighten{F}}$
by lemma 
\lref{tighteningCommutesWithConcatenatingAnIdempotent}%zzzzz
(i)
and the first element of
$\tighten{\tuple{\rng{i}} \cnc \tighten{F}}$
is $\rng{i} \circ t_1$
by lemma \lref{alpha} (i)
and so we have
\begin{align*}
\rst{\tighten{\tuple{\rng{i}} \cnc F}}
&= \tuple{\rst{\rng{i} \circ t_1}}
          && \mbox{by definition of the support operator for zigzags,} \\
&= \tuple{\rng{i} \circ \rst{t_1}}
          && \mbox{by lemma \lref{bisupportLemma} (i),} \\
&= \tuple{\rng{i}} \circ \tuple{\rst{t_1}}
          && \mbox{\parbox[t]{5cm}{by definition of composition and since singleton zigzags are tight by definition}} \\
&= \tuple{\rng{i}} \circ \rst{\tighten{F}}
          && \mbox{by definition of the support operator.}
\end{align*}
\item %(ii)
Follows by duality since it is the dual of (i).
\end{enumerate}
\end{proof}

% lemma howToTightenAConcatenation
\newcommand{\tighteningLeft}{\Tighten{ F \cnc 
                                  \rst{\tighten{G}}
                                      }
                            }

\newcommand {\tighteningRight}{\Tighten{ 
                                      \rng{\tighten{F}} 
                                  \cnc G
                                        }
                              }
\newcommand {\theTightening}{\tighteningLeft \cnc \tighteningRight}


\begin{lemma}
\llabel{howToTightenAConcatenation}
If $F$ and $G$ are zigzags that compose so that
$F \cnc G$ is defined then
\begin{align*}
\tighten{F \cnc G} 
    &= \theTightening
\end{align*}
\end{lemma}
\begin{proof}
We need to show two things. First that  $\theTightening$
is a tight zigzag and secondly that for any 
zigzag $H$, if $H \leq F \cnc G$ then $H \leq \theTightening$.

For the first part we know (lemma \lref{bisupportLemma} (iv)) 
that if  $f$ and $g$ are morphisms 
of the bisupport category \catcw 
such that $f \circ g$ defined and $\rng{f} = \rst{g}$ 
then $\rst{f \circ g}=\rst{f}$ and 
 $\rng{f \circ g} = \rng{g}$.
As a consequence,
 for any tight zigzags
$Z_1$ and $Z_2$ such that the concatenation $Z_1 \cnc Z_2$ is defined then if  $\rng{Z_1} = \rst{Z_2}$ then
$Z_1 \cnc Z_2$ is tight. 

It suffices, therefore to show that $\rng{Z_1} = \rst{Z_2}$
when $Z_1$ is $\tighteningLeft$ and
$Z_2$ is $\tighteningRight$
but this is the case because by lemma \lref{beta} both left and right hand side evaluate  to 
$\rng{\tighten{F}} \circ \rst{\tighten{G}}$.\highlight{CHECK}

To complete the proof we assume that $H$ is a tight zigzag such that $H \leq F \cnc G$ and we have to prove that 
$H \leq \theTightening$. For this we let
$F=\zigzagTuple{f}{n}$, $G=\zigzagTuple{g}{m}$
and we let $H=\zigzagLeadingTuple{h}{n+m-1}$.

Note that we have $h_n \leq f_n \circ g_1$ and therefore
we are able to
use construction \lref{fghConstruction}. 
We shall make use of $h_L$ and $h_R$ as defined there and translating into the situation here, they are given by
\[
h_L = \rst{h_n} \circ f_n \circ \rst{g_1 \circ \rng{h_n}}
\]
and  
\newcommand{\hR}{\rng{\rst{h_n} \circ f_n} \circ g_1 \circ \rng{h_n}}
\[
h_R = \hR.
\]
and, as described in that construction, are such that
\begin{equation}
\label{rstOfhL}
\rst{h_L}=\rst{h_n}
\end{equation}
\begin{equation}
\label{rngOfhR}
\rng{h_R}=\rng{h_n}.
\end{equation}

Now,  
\begin{equation}
\label{inequalityOneLeft}
\zigzagLeftFactorTupleLedOut{h}{n}{f_n} \leq F
\end{equation}
is the case because of the component inequalities all but the final one hold because $H \leq F \cnc G$
and the final one, $\rst{h_n} \circ f_n \leq f_n$, holds because
the left hand side is a support idempotent composed 
with the right hand side (lemma \lref{idempotentCompositionAndOrdering} (i)).

Likewise
\begin{equation}
\label{inequalityOneRight}
\zigzagRightFactorTailTupleLedIn{h}{n}{m}{g_1} \leq G
\end{equation}
is the case because of the component inequalities: all but the initial one hold because
$H \leq F \cnc G$,
and the initial one,
\[
g_1 \circ \rng{h_n} \leq g_1,
\]
holds because the left hand side is the right hand side composed with a cosupport idempotent
(lemma \lref{idempotentCompositionAndOrdering} (ii)).

Because $h_n \leq f_n \circ g_1$ and the support operator is monotone
then we have
$\rst{h_n} \leq \rst{f_n \circ g_1}$
and because $\rst{f_n \circ g_1} \leq \rst{f_n}$ (by lemma \lref{idempotentCompositionAndOrdering}(v)), 
so transitively we arrive at $\rst{h_n} \leq \rst{f_n}$, or, 
equivalently, $\rst{h_n} \circ \rst{f_n} = \rst{h_n}$, which then yields $\rst{\rst{h_n} \circ f_n} = \rst{h_n}$, by use of wR.4.
This means  that the zigzag  
$\zigzagLeftFactorTupleLedOut{h}{n}{f_n}$ 
is tight which implies that
 from (\ref{inequalityOneLeft})  we can deduce that
\[
\zigzagLeftFactorTupleLedOut{h}{n}{f_n} \leq \tighten{F}
\]
and from this, by the way, that
\[
\rng{\zigzagLeftFactorTupleLedOut{h}{n}{f_n}} 
\leq \rng{\tighten{F}}
\]
and thereby that
\begin{equation}
\label{inequalityThreeLeft}
\tuple{\rng{\rst{h_n} \circ f_n}} \leq \rng{\tighten{F}}
\end{equation}

Now we show that the zigzag 
$\zigzagRightFactorTailTupleLedIn{h}{n}{m}{g_1}$,
which is the right hand side of (\ref{inequalityOneRight}),
is tight. It is evident that its joins except for its initial join are tight because $H$ is tight.
In the case of the initial join,
given that the tightness of $H$ gives us 
that $\rng{h'_n}=\rng{h_n}$,
the tightness condition can be written as
$\rng{g_1 \circ \rng{h_n}} = \rng{h_n}$.
We can show this as follows.  
From the fact that
$h_n \leq f_n \circ g_1$ we get
$\rng{h_n} \leq \rng{f_n \circ g_1}$,
and because generally $\rng{f_n \circ g_1} \leq \rng{g_1}$ (by lemma \lref{idempotentCompositionAndOrdering}(vi)),
transitively we get 
$\rng{h_n} \leq \rng{g_1}$
which, by definition, is $\rst{\rng{h_n}} \circ \rng{g_1} = \rng{h_n}$
which simplifies and rearranges to
$\rng{g_1} \circ \rng{h_n} = \rng{h_n}$.
This is also a simplification of
$\rng{g_1 \circ \rng{h_n}} = \rng{h_n}$, as required for tightness.

Having shown that its right hand side is tight
we can deduce from (\ref{inequalityOneRight}) that
\[
\zigzagRightFactorTailTupleLedIn{h}{n}{m}{g_1} \leq \tighten{G}
\]
and thus 
\begin{equation}
\label{inequalityThreeRight}
\tuple{\rst{ g_1 \circ \rng{h_n}}} \leq \rst{\tighten{G}}
\end{equation}

Putting (\ref{inequalityOneLeft})
 and (\ref{inequalityThreeRight}) together (using the fact that
 composition in \catcw is monotone) we get 
 
\[
\zigzagLeftFactorTupleLedOut{h}{n}{f_n} 
          \cnc \tuple{\rst{ g_1 \circ \rng{h_n}}}
              \leq F \cnc \rst{\tighten{G}}
\]
which is to say
\[
\zigzagLeftFactorTupleLedOut{h}{n}
        {f_n \circ \tuple{\rst{ g_1 \circ \rng{h_n}}}}
              \leq F \cnc \rst{\tighten{G}}
\]
which is 
\begin{equation}
\label{inequalityFourLeftRewrittenAgain}
\zigzagTupleReplaceOut{h}{n}
        {h_L}
              \leq F \cnc \rst{\tighten{G}}
\end{equation}
and since the left hand side of this inequality is tight
(by (\ref{rstOfhL}) and since $\rst{h_n} = \rst{h'_{n-1}}$
 as $H$ is tight)
then we get
\begin{equation}
\label{inequalityFiveLeft}
\zigzagTupleReplaceOut{h}{n}
        {h_L}
              \leq \Tighten{F \cnc \rst{\tighten{G}}}.
\end{equation}

This supports the left hand side of our target inequality, for the right hand side we
put (\ref{inequalityOneRight}) 
and (\ref{inequalityThreeLeft}) together to get
\[
\tuple{\rng{\rst{h_n} \circ f_n}} \cnc \zigzagRightFactorTailTupleLedIn{h}{n}{m}{g_1} 
       \leq \tuple{\rng{\tighten{F}}} \cnc G
\]
which is to say
\[
\zigzagTailTupleReplaceIn{h}{n}{m}{\hR} 
       \leq \tuple{\rng{\tighten{F}}} \cnc G 
\]
i.e.
\[
\zigzagTailTupleReplaceIn{h}{n}{m}{h_R} 
        \leq  \tuple{\rng{\tighten{F}}} \cnc G 
\]
and since the left hand side of this inequality is tight
(by (\ref{rngOfhR}) and since $\rng{h_n} = \rng{h'_n}$
 as $H$ is tight)  we deduce
\begin{equation}
\label{inequalityFiveRight}
\zigzagTailTupleReplaceIn{h}{n}{m}
                {h_R} 
                 \leq  \Tighten{ \rng{\tighten{F}} \circ G}
\end{equation}

Finally, putting (\ref{inequalityFiveLeft}) 
and (\ref{inequalityFiveRight}) together
(using the fact that
 composition in \catcw is monotone)
we can  simplify as  $h_L \circ h_R = h_n$ and we get 
\[
H \leq \theTightening,
\]
as required to complete the proof.
\end{proof}
 % end of lemma howToTightenAConcatenation

\newcommand{\tightCircTightLeft}[2]{\tighten{#1 \circ \rst{#2}}}
\newcommand{\tightCircTightRight}[2]{\tighten{\rng{#1} \circ #2}}
\newcommand{\tightCircTight}[2]
       {\tightCircTightLeft{#1}{#2} \cnc \tightCircTightRight{#1}{#2}}

\newcommand{\FcircGleft}{\tightCircTightLeft{F}{G}}
\newcommand{\FcircGright}{\tightCircTightRight{F}{G}}
\newcommand{\FcircG}{\tightCircTight{F}{G}}

\begin{corollary}
\llabel{howToTightenAConcatenationOfTights}
If $F$ and $G$ are tight zigzags such that
$F \cnc G$ is defined then
\[
\tighten{F \cnc G} = \FcircG
\]
\end{corollary}

\begin{lemma}
\llabel{tighteningCommutesWithConcatenation}
If $F$ and $G$ are zigzags for which the concatenation $F\cnc G$ is defined
then 
$ T(F \cnc  G) = T(T(F) \cnc T(G)) $
\end{lemma}


\newcommand{\tightenedComposite}[1]
{
\zigzagOneAndAHalf{w}{x}{#1}
\begin{array}{c}
\dots \\[0.8cm]
\dots
\end{array}
\begin{array}{c c}
\Ob{x}{n} \\[0.8cm]
& \Ob{w}{n}=
\end{array}
\begin{arrows}
\ncarr{xn}{wn} \blabel{#1'_{n-1}}[0.2][0.4] %0.2 was 0.45
\end{arrows}
\kern-15pt
\arraycolsep=7pt
\zigzagFlex{y}{z}{#1}{n}{n+m}
}

\begin{proof}
\begin{align*}
\tighten{F \cnc G}
 &= \tighten{F \cnc \rst{\tighten{G}}} 
                    \cnc \tighten{\rng{\tighten{F}}\cnc G}
    &&\mbox{by lemma \lref{howToTightenAConcatenation},} \\
 &= \tighten{\tighten{F} \cnc \rst{\tighten{G}}} 
                    \cnc \tighten{\rng{\tighten{F}}\cnc \tighten{G}}
    &&\mbox{by lemma \lref{tighteningCommutesWithConcatenatingAnIdempotent},}\\
 &= \tighten{\tighten{F} \cnc \tighten{G}}
    &&\mbox{by lemma \lref{howToTightenAConcatenationOfTights}.}
\end{align*} 

\end{proof}

\subsection{The category of tight zigzags}
\llabel{categorytightzigzags}
If \catcw is a range category then we can define the category
$Z_C$ to be the category whose objects are the objects of \catcw
and such that $Hom_{Z_C}(a,b)$ is the set of tight zigzags of any length $n$
$$
\arraycolsep=7pt
\zigzag{x}{y}{f}{n}
$$
such that $x_1=a$ and $y_n=b$.


Composition in $Z_C$ is defined by composing the zigzags and then tightening  thus
\[
F \circ G =_{def} T(F :: G).
\]
Composition is associative since 
\begin{align*}
(F \circ G) \circ H 
                &= T(T(F \cnc G) \cnc H) &&   \mbox{by definition,} \\
                &= T(T(F \cnc G) \cnc T(H))
                        &&\mbox{since $H$ is tight,}\\
                &= T((F \cnc G) \cnc H)  
                        &&\mbox{by lemma \lref{tighteningCommutesWithConcatenation},}\\
                &= T(F \cnc (G \cnc H))
                        &&\mbox {since $\cnc$ is associative,} \\
                &= T(T(F) \cnc T(G \cnc H))       
                        &&\mbox{by lemma \lref{tighteningCommutesWithConcatenation},}\\
                &= T(F \cnc T(G \cnc H))))
                        &&\mbox{since $F$ is tight,} \\
                &= F \circ (G \circ H) &&   \mbox{by definition.}                 
\end{align*}

Identity morphisms in $Z_C$ consist of the identity morphisms of \catcw considered 
as zigzags on length 1.

\begin{worktt}
\subsection{The category of tight zigzags $Z_C$ is a category with support}

If $F=\zigzagTuple{f}{n}$ is a morphism in $Z_C$ then define $\rst{F}$ to be $\tuple{\rst{f_1}}$.

We need show

R.1 For $F:a \morph b$ in $Z_C$ $$\rst{F} \circ F =F$$.

R.2. If \FGsourcediag in $Z_C$ then
$$\rst{G} \circ \rst{F}=\rst{F} \circ \rst{G}.$$

R.3. If \FGsourcediag in $Z_C$ then
$$\rst{\rst{F} \circ G} = \rst{F} \circ \rst{G}$$.

wR.4. If $\sequentialdiag{a}{b}{c}{F}{G}$ in $Z_C$ then
\[
\rst{F \circ \rst{G}} = \rst{F \circ G}.
\]

R.1 follows because in \catcw, $f_0 \circ \rng{f_0} = f_0$. 

R.2 follows because in \catcw $\rng{f_0} \circ \rng{g_0} = \rng{g_0} \circ \rng{f_0}$.

R.3 follows because
\begin{align*}
\rst{\rst{F} \circ G} 
                   &= \rst{T(\tuple{\rst{f_0} \circ g_0, g'_0,...g_m})}\\
                   &= \rst{\tuple{\rst{f_0} \circ g_0,s'_0,...s_m}}
                                             &&\mbox{\parbox{5cm}{for some $s'_0,...s_m$
                                                       by some new lemma about tightening when only head position loose,}}\\
                   &=\tuple{\rst{\rst{f_0} \circ g_0}} 
                                            && \mbox{by definition}\\
                   &=\tuple{\rst{f_0} \circ \rst{g_0}}  
                                            && \mbox{by R.3}\\
                   &=\rst{F} \circ \rst{G}      
\end{align*}

If $F$ is of length 1 then wR.4 holds because
it holds in the underlying range category \catcw. 
If $F$ is of length $n > 1$ then wR.4 follows because lemma
\lref{midpointLemma} informs us that $\rst{f \circ G}$ depends only on the $\rst{g_1}$ component of $G$ i.e. only on $\rst{G}$.
\end{worktt}

\begin{worktt}
\workt{The K.4 part needs moving after contraction}
\begin{lemma}
\llabel{compatibleConverseLemma}
Defining  the converse of a tight zigzag $\zigzagTuple{f}{n}$ 
to be the zigzag $\tuple{\rng{f_n},f_n,f'_{n-1},...f'_1,f_1,\rst{f_1}}$
\commentary{contract the converse if necessary}
equips the category of \highlight{contracted} tight zigzags over a bisupport category with a compatible converse.
\end{lemma}
\begin{proof}
We have to prove K.1, K.2 and K.3 and \highlight{K.4}.

K.1 is trivial because when converse is taken and then taken again then the result can be \highlight{contracted} to the given zigzag. \commentary{draw this}


\begin{aside}
This below doesn't quite add up --- feel we may need contrcation somewhere. Unless this has dissapeared in side of an argument regarding converse of a converse.
\end{aside}
K.2 \highlight{K.4 ??} follows because
\begin{align*} 
\converse{(F \circ G)}  
&= \converse{\biggl(\FcircG \biggr)}
&&\mbox{by corollary \lref{howToTightenAConcatenationOfTights},} \\
&= \converse{\biggl(\FcircGright \biggr)} \cnc \converse{\biggl(\FcircGleft\biggr)} 
&&\mbox{because $\converse{(Z_1 \cnc Z_2)}=\converse{Z_2} \cnc \converse{Z_1}$,\highlight{really? I don't think so}} \\
&=
\tighten{\converse{\bigl( 
                      \rng{F} \cnc G
                    \bigr)}  
        } 
              \cnc 
\tighten{\converse{\bigl(
                       F \cnc \rst{G}
                     \bigr)}
        }
&&\mbox{\parbox[t]{6cm}{  since $\tighten{\converse{Z}}=\converse{\tighten{Z}}$,
                        }}\\
&=
\tighten{\converse{G} \cnc \rng{F} 
        } 
              \cnc 
\tighten{\rst{G} \cnc \converse{F}
        }
&&\mbox{\parbox[t]{6cm}{  using $\converse{(Z_1 \cnc Z_2)}=\converse{Z_2} \cnc \converse{Z_1}$ and since a singleton is unchanged when a converse is taken,
                        }} \\
&=
\tighten{\converse{G} \cnc \rst{\converse{F}} 
        } 
              \cnc 
\tighten{\rng{\converse{G}} \cnc \converse{F}
        }
&&\mbox{\parbox[t]{6cm}{ since $\rst{\converse{Z}}=\rng{Z}$ and, dually,
                                $\rng{\converse{Z}} = \rst{Z}$ ,
                        }} \\
&= \converse{G} \circ \converse{F}
&&\mbox{by corollary \lref{howToTightenAConcatenationOfTights}.}
\end{align*}

K.3 follows simply because
\begin{align*}
\rst{\converse{F}}&= \rng{f_n} 
                  && \mbox{from definitions of converse and
                           of \rstItself for a tight zigzag,} \\
                  &= \rng{F}
                  && \mbox{from definition of \rngItself 
                              for a tight zigzag.}
\end{align*} 
\end{proof}
\end{worktt}

\newcommand{\BiSup}{BiSup}
\newcommand{\BiSupConv}{\converse{\BiSup}}

We can define a functor $CZ: \BiSup \morph \BiSupConv$ that
is left adjoint  to the forgetful functor $U: \BiSupConv \morph \BiSup$
by defining, $CZ(B)$ to be the category of contracted tight zigzags over $B$
for any bisupport category $B$ 
and by defining $CZ(F)$, for any bisuport functor $F: B \morph B'$, 
as pointwise application of the functor $F$ 
to the elements of the zigzags of $CZ(B)$ 
followed by contraction 
(necessary because $F$ may map a non-idempotent to an idempotent). 
Because as a bisupport functor $F$ preserves \rstItself and \rngItself 
then the image under $F$ of tight zigzags are tight.

We need show that $CZ(F)(\converse{Z})=\converse{CZ(F)(Z)}$. 
This holds because $CZ(F)$ is defined pointwise and because contraction
commutes with converse.

\newcommand{\funU}{\cat{U}}
\newcommand{\funUw}{\cat{U}\ }
\newcommand{\funCZ}{\cat{CZ}}
\newcommand{\funCZw}{\cat{CZ}\ }

\begin{lemma}
If \catcw is a bisupport category then
the category of contracted tight zigzags of \catcw
is the free bisupport category with compatible converse over \catcw.
so that $\funU(\catb)$ is the underlying bisupport category of \catcw
\end{lemma}
\begin{proof}
Define $\eta_\catc: \catb \morph \funU(\funCZ(\catb))$ be 
defined as the mapping that maps morphisms of \catbw to singletonn zigzags.
This is a bisupport functor (Because?). 
This will be the unit of the adjunction.
We  show that if $\gamma: \catb \morph \funU(\catc)$ is a bisupport functor
from a  bisupport category \catbw to the underlying category
$\funU(\catc)$ of a  bisupport category  with converse \catcw
then there is a unique converse preserving bisupport functor
$\psi(\gamma):\funCZ(\catb) \morph \catc$ such that 
\[
\upTriangleShape{\funU(\catc)}{\catb}{\funU(\funCZ(\catb))}
     %{top}[height]{left}[lpadding]{right}[rpadding]
\begin{arrows}
% \drawMorphism
%      {src}{dest}{labelside(a|b)}{label}[npos]
\drawMorphism{left}{top}{a}{\gamma}
\drawMorphism{right}{top}{b}{\funU(\psi(\gamma))}
\drawMorphism{left}{right}{b}{\eta_\catc}
\end{arrows}
\]
commutes.

The objects of the zigzag category are the objects of the underlying category \catbw.
Define $\psi(\gamma)$ to map these objects of the zigzag category
as $\gamma$ maps the corresponding objects of \catbw:
\begin{equation}
\label{adjunctionTriangle}
\begin{array}{rcl}
a &\overset{\psi(\gamma)}{\mapsto}& \gamma(a),\\
\end{array}
\end{equation}

The morphisms are contracted zigzags. 
Define $\psi(\gamma)$ on a typical morphism,  
a  zigzag $\zigzagTuple{f}{n}$, as follows:
\[
\begin{array}{rcl}
\zigzagTuple{f}{n}
&\overset{\psi(\gamma)}{\mapsto}&
\gamma(f_1)
\circ \converse{(\gamma(f'_1))}
\circ \gamma(f_2)
\circ \cdots
\circ \converse{(\gamma(f'_{n-1}))}
\circ \gamma(f_n).
\end{array}
\]

We have to show that $\psi(\gamma)$ is a functor and that it preserves
\rstItself, \rngItself and \converseItself.

Preserves \rstItself.

Preserves \rngItself.

Because $\psi(\gamma)$ preserves \rstItself and \rngItself and because
tightening uses exclusively  \rstItself and \rngItself it follows that
$\psi(\gamma)$ preserves tightening.

Also note that from its definition 
it is easily seen that $\psi(\gamma)$  preserves $\cnc$.

Now functoriality of $\psi(\gamma)$ follows because
$F \circ G = \FcircG$ and therefore $\psi(\gamma)(F \circ G)
=\tightCircTight{{\psi(\gamma)(F)}}{\psi(\gamma)(G)}
= \psi(\gamma)(F)\circ \psi(\gamma)(G)$

Finally see that $\psi(\gamma)$ preserves \converseItself 
by using the bisupport axioms in the target bisuppport category $\funU(\catc)$.

Finally we have to show that $\psi(\gamma)$, as defined, is the unique 
converse preserving bisupport functor from $\funCZ(\catb)$ to \catcw 
such that diagram (\ref{adjunctionTriangle}) commutes. This follows
because every morphism in $\funCZ(\catb)$ is a composition 
of singleton zigzags and their converses i.e. of morphisms that are either
in the image of $\eta_\catc$ or are converses of such. Since the mapping of these singletons must agree with the mapping of $\gamma$ then the mapping of
$\psi(\gamma)$, as we have defined it, is unique as composition and \converseItself are preserved.  

\end{proof}

\subsection{Tightening Length Two Zigzags}
{%howToTighteningaZigZagOfLengthTwo

\begin{lemma}
\llabel{howToTighteningaZigZagOfLengthTwo}
If $Z=\tuple{z_1,z'_1,z_2}$ is a zigzag of length $2$
then

\[
\tighten{Z}=\tuple{ \twoTighteningOne{z},
                     \,\twoTighteningTwo{z},
                     \,\twoTighteningThree{z}}
\]
\end{lemma}
\begin{proof}
\newcommand {\longTighteningTwo}{\rst{z_2} \circ z'_1 \circ \rng{t}}
\newcommand {\longTighteningThree}{\rst{z'_1 \circ \rng{t}} \circ z_2}
Use lemma \lref{leftToRightTighteningSubLemma} with $n=2$. The precursor zigzag of length $n-1$ is simply the singleton zigzag \nudgeup{0.55cm}$\tuple{\twoTighteningOne{z}}$ and this being a  singleton is tight and therefore its own tightening. Therefore the tightening of $Z$ is 
\nudgeup{0.55cm}$\tuple{t, \longTighteningTwo,\longTighteningThree}$
where $t = \twoTighteningOne{z}$. 

Now note that 
\begin{equation}
\label{tighteningTwoIntermediate}
z'_1 \circ \rng{t} = z'_1 \circ \rng{\rst{z_2} \circ z'_1} \circ \rng{z_1}
\end{equation}  
because
\begin{align*}
z'_1 \circ \rng{t} 
     &= z'_1 \circ \rng{\twoTighteningOne{z}} 
     && \mbox{substituting $t$,}\\
     &= z'_1 \circ \rng{z_1} \circ \rng{\rst{z_2} \circ z'_1} 
     && \mbox{by  axiom C.3,}\\
     &= z'_1 \circ \rng{\rst{z_2} \circ z'_1}  \circ \rng{z_1}
     && \mbox{by axiom C.2}
\end{align*}
and using this we can simplify
$\longTighteningTwo$  to $\twoTighteningTwo{z}$ as follows
\begin{align*}
\longTighteningTwo 
    &= \rst{z_2} \circ z'_1 \circ \rng{\rst{z_2} \circ z'_1} \circ \rng{z_1}
                   && \mbox{by (\ref{tighteningTwoIntermediate}),}\\
    &= \twoTighteningTwo{z}
                   && \mbox{by axiom C.1,}
\end{align*}
and  $\longTighteningThree$ simplifies to $\twoTighteningThree{z}$ as follows
\begin{align*}
\longTighteningThree 
     &= \rst{z'_1 \circ \rng{\rst{z_2} \circ z'_1} \circ\rng{z_1}} \circ z_2
     && \mbox{by (\ref{tighteningTwoIntermediate}),}\\
     &= \rst{\rst{z_2} \circ z'_1 \circ \rng{\rst{z_2} \circ z'_1} \circ\rng{z_1}
      } \circ z_2
     && \mbox{by  lemma \lref{supportLemma} (v)}\\
     &= \rst{\rst{z_2} \circ z'_1 \circ\rng{z_1}} \circ z_2 
     && \mbox {by axiom C.1,}\\ 
     &= \twoTighteningThree{z}
     && \mbox{by a second use of lemma \lref{supportLemma} (v).}
\end{align*} 
Hence result.
\end{proof}

\subsection{How to Tighten ThreeFold Concatenations}
\begin{lemma}[howToTightenAThreefoldConcatenation]
\llabel{howToTightenAThreefoldConcatenation}
If $L$, $Z$ and $R$ are zigzags that compose so that
$L \cnc Z \cnc R$ is defined then if we let
$P_1 =\tighten{ L \cnc \abrst{{Z \cnc r}}}$
$P_2 =\Tighten{l\cnc Z \cnc r}$,
and   $P_3 = \tighten{ \abrng{l \cnc Z} \cnc R}$,
where $l = \abrng{L}$
and   $r = \abrst{R}$,
then
\[
\Tighten{L \cnc Z \cnc R} = 
     P_1 \cnc P_2 \cnc  P_3.
\]
\end{lemma}
\begin{proof}
First note that
\begin{align*}
\tighten{ L \cnc Z \cnc R} 
               &=\tighten{ (L \cnc Z) \cnc R} \\ 
               &= \tighten{ (L \cnc Z) \cnc \abrst{R}} 
                     \cnc
                  \tighten{ \abrng{L \cnc Z} \cnc R} 
               &&\mbox{by lemma 
                    \lref{howToTightenAConcatenation}}
\end{align*}
and then that 
\begin{align*}
\tighten{ (L \cnc Z) \cnc \abrst{R}}  
        &= \tighten{ L \cnc (Z \cnc \abrst{R})}        \\
        &= \tighten{ L \cnc \abrst{{Z \cnc \abrst{R}}}} 
              \cnc
           \tighten{ \abrng{L} \cnc (Z \cnc \abrst{R})} 
           &&\mbox{by lemma 
                    \lref{howToTightenAConcatenation},} \\
        &= P_1 \cnc P_2 
           &&\mbox{where $P_1$, $P_2$ as given above}       
\end{align*}   
and also 
\begin{align*}
\tighten{ \abrng{L \cnc Z} \cnc R}
     &= \tighten{ \abrng{\abrng{L} \cnc Z} \cnc R}
          &&\mbox{
               \parbox[t]{6cm}
               {
               Because it follows from lemma \lref{howToTightenAConcatenation}
               that for all composable zigzags $F$ and $G$, 
                $\abrng{F \cnc G} = \abrng{\abrng{F} \cnc G}$
               }
          }\\
     &= P_3 &&\mbox{where $P_3$ as given above.}
\end{align*}
\end{proof}

